\clearpage
\VerdierPhysicalPageAnchor{0150}%
\begingroup
\small
\setlength{\parskip}{0pt}%
\setlength{\emergencystretch}{1.5em}%

\VerdierSemanticAnchor{ch2-2-3-5-w-continuation}
\noindent et comme $B$ est un objet de $\mathcal B$ et $s$ un élément de
$S_{\mathcal D}(\mathcal A\cap\mathcal B)$, l'objet $W$ est isomorphe à un
objet de $\mathcal B$ (1.2.4).

\VerdierSemanticAnchor{ch2-2-3-5-i-implies-ii-conclusion}
Par suite, d'après la première assertion, le morphisme $f$ se factorise à
travers un objet de $\mathcal A\cap\mathcal B$. Le morphisme $Q_3(f)$ est
donc nul, et par suite $m=0$.

\VerdierSemanticAnchor{ch2-2-3-5-ii-implies-i}
Démontrons maintenant l'implication ii) $\Rightarrow$ i). Il suffit de
montrer, d'après a), que tout morphisme dans $\mathcal D$,
$g\colon B\longrightarrow X$, où $B$ est un objet de $\mathcal B$ et $X$
un objet de $\mathcal A$, se factorise à travers un objet de
$\mathcal A\cap\mathcal B$.

\VerdierSemanticAnchor{ch2-2-3-5-q3-g-zero}
On sait par hypothèse que
$Q_3(g)\in
\operatorname{Hom}_{\mathcal D/(\mathcal A\cap\mathcal B)}(B,X)$ est nul.

\VerdierSemanticAnchor{ch2-2-3-5-t-morphism}
Par suite (2.2.4), il existe un morphisme
$t\colon X\longrightarrow X'$, $t\in
S_{\mathcal D}(\mathcal A\cap\mathcal B)$, tel que $tg=0$.

\VerdierSemanticAnchor{ch2-2-3-5-distinguished-triangle}
Par définition de l'ensemble
$S_{\mathcal D}(\mathcal A\cap\mathcal B)$, il existe un triangle distingué
de $\mathcal D$ de la forme $(B',X,X',i,t,w)$, où $B'$ est un objet de
$\mathcal A\cap\mathcal B$.

\VerdierSemanticAnchor{ch2-2-3-5-proof-complete}
D'après (1.2.1), l'égalité $tg=0$ implique que le morphisme $g$ se
factorise à travers $B'$, ce qu'il fallait démontrer. Ceci achève la
démonstration de la proposition (2.3.5).

\VerdierSemanticAnchor{ch2-2-3-5-left-localizing-definition}
\medskip
On dit qu'une sous-catégorie triangulée pleine $\mathcal A$ de $\mathcal D$
est \emph{$\mathcal B$-localisante à gauche} si la sous-catégorie
$\mathcal A^\circ$ de $\mathcal D^\circ$ est $\mathcal B^\circ$-localisante
à droite.
\endgroup
